\section{习题}

\subsection{基础过关题}

\begin{enumerate}

\item \textbf{（单项选择题）} 在 UNIX 类文件系统中，目录项（directory entry）与 i-node 的分工是
\begin{center}
\begin{tabular}{p{0.47\linewidth}p{0.47\linewidth}}
(A) 目录项保存文件名和文件的全部属性，i-node 只保存数据块号 & (B) 目录项保存文件名和 i-node 号，文件类型、长度、物理块号等属性保存在 i-node 中 \\[4pt]
(C) 一个 i-node 只能被一个目录项引用，所以 i-node 中的链接计数恒为 1 & (D) 删除文件时只要删掉目录项，其 i-node 和数据块就立即被回收
\end{tabular}
\end{center}

\ifshowsolutions\par\noindent\textbf{解析：} 思路：先看“按名存取”需要什么，再看“文件属性”放在哪里。目录的职责是把用户给的文件名翻译成文件实体，所以目录项至少要有文件名和指向文件实体的编号；UNIX 把除文件名以外的属性全部放进 i-node，目录项只剩“名字 + i-node 号”。

由此 (A) 错：目录项不存全部属性。(C) 错：多个目录项可以指向同一 i-node，这正是硬链接，链接计数记录的就是目录项个数。(D) 错：删除目录项只是把链接计数减 1，计数减到 0 才释放 i-node 和数据块，否则其它硬链接会失效。\boxed{\textbf{答案：}(B)}\else \answerarea \fi

\vspace{8pt}

\item \textbf{（单项选择题）} 关于文件的物理结构（分配方式），下列说法正确的是
\begin{center}
\begin{tabular}{p{0.47\linewidth}p{0.47\linewidth}}
(A) 连续分配支持随机访问，但文件增长时若后继盘块已被占用，可能要整体搬移 & (B) 隐式链接分配中，访问文件的第 $i$ 个逻辑块只需读一次盘 \\[4pt]
(C) 索引分配中，索引块里存放的是文件的数据内容 & (D) 连续分配在多次分配与回收之后不会产生外部碎片
\end{tabular}
\end{center}

\ifshowsolutions\par\noindent\textbf{解析：} 连续分配用“起始块号 + 长度”描述文件，第 $i$ 块地址可直接算出，所以随机访问快；代价是必须占用连续空间，文件变长时若后面紧邻的块已被别的文件占用，就只能把整个文件搬到足够大的空闲区，并修改目录项中的起始块号，因此 (A) 正确。

(B) 错：隐式链接的指针存在每个数据块里，必须从第一块起逐块读，访问第 $i$ 块一般需要 $i$ 次读盘。(C) 错：索引块存的是数据块号，不是数据本身。(D) 错：连续分配正是外部碎片的典型来源——空闲区被已分配区切成许多不连续的小块，总空闲量够但凑不出连续空间。\boxed{\textbf{答案：}(A)}\else \answerarea \fi

\vspace{8pt}

\item \textbf{（填空题）} 某系统用位示图管理磁盘空闲空间：磁盘共有 $2^{26}$ 个盘块，位示图每个字 64 位，约定“1 表示已分配、0 表示空闲”，盘块号、字号、位号均从 0 开始编号，位示图第 $i$ 个字的第 $j$ 位（自低位向高位计）对应盘块号 $b=i\times 64+j$。
（1）位示图共占用 \underline{\hspace{2.6cm}} 字节；
（2）盘块号 $2\,000\,003$ 对应的字号 $i=$ \underline{\hspace{1.8cm}}，位号 $j=$ \underline{\hspace{1.4cm}}。

\ifshowsolutions\par\noindent\textbf{解析：} 思路：位示图用 1 位标记 1 个盘块，位示图总位数就等于盘块总数，先把“位”折成“字”，再折成“字节”。

（1）总位数为 $2^{26}$，每字 64 位，故字数 $=2^{26}/2^{6}=2^{20}=1\,048\,576$ 个。每个字 8 字节，位示图占 $2^{20}\times 8=2^{23}=8\,388\,608$ 字节，即 8 MiB。

（2）$2\,000\,003 \div 64 = 31\,250$ 余 3，因为 $31\,250\times 64 = 2\,000\,000$。所以字号 $i=31\,250$，位号 $j=3$，该位对应的盘块号正是 $31\,250\times 64+3=2\,000\,003$。

\boxed{\text{答案：（1）$8\,388\,608$ 字节（8 MiB）；（2）$i=31\,250$，$j=3$}}\else \answerarea[5] \fi

\vspace{8pt}

\item \textbf{（单项选择题）} 关于磁盘调度算法 SCAN 与 C-SCAN，下列说法正确的是
\begin{center}
\begin{tabular}{p{0.47\linewidth}p{0.47\linewidth}}
(A) SCAN 对靠近磁盘两端的柱面比对中间柱面更公平 & (B) C-SCAN 只在一个方向服务请求，返程快速移回起始端且不服务，因此各柱面的等待时间更均匀 \\[4pt]
(C) 对同一请求序列，C-SCAN 的总寻道长度一定小于 SCAN & (D) LOOK 与 SCAN 的唯一区别是 LOOK 必须移到磁盘最外端才折返
\end{tabular}
\end{center}

\ifshowsolutions\par\noindent\textbf{解析：} SCAN 像电梯一样来回扫描，两端柱面刚被扫过就要等磁头走完一个来回，等待时间反而比中间柱面长，故 (A) 说反了。

C-SCAN 只在磁头前进方向服务，扫到端点后直接快速返回起点、返程不服务，于是每个柱面的等待时间接近“一个方向上的扫描周期”，分布更均匀，(B) 正确。

(C) 错：返程那一段移动量 C-SCAN 也要花时间（有的实现还把它计入寻道），总寻道长度常常更大，本题问的是“一定小于”，不成立。(D) 错：LOOK 的特点恰恰是看到该方向上没有更远的请求就折返，不必到达物理端点。\boxed{\textbf{答案：}(B)}\else \answerarea \fi

\vspace{8pt}

\item \textbf{（填空题）} 某文件系统盘块大小为 4 KiB，每个盘块号占 4 字节，按 $1\text{ KiB}=2^{10}$ 字节计。
（1）只采用一级索引时，一个索引块能描述的最大文件长度为 \underline{\hspace{2.4cm}}；
（2）若改用 FAT16（FAT 表项 16 位，按 $2^{16}$ 个表项计）管理同样 4 KiB 的盘块，则一个分区最大为 \underline{\hspace{2.4cm}}。

\ifshowsolutions\par\noindent\textbf{解析：} 思路：索引块能装多少个盘块号，取决于“块大小 $\div$ 每个块号占的字节数”；一个块号对应一个数据块，所以再乘块大小就是能描述的文件长度。FAT 同理：表项位数决定块号个数上限，块号个数乘块大小就是分区上限。

（1）一个索引块可存 $4\text{ KiB}/4\text{ B}=4096/4=1024$ 个块号，每个块号指向 4 KiB 数据，故最大文件长度 $=1024\times 4\text{ KiB}=4\text{ MiB}$（此时索引块自身占 1 块，通常不计入文件数据长度）。

（2）16 位表项能表示 $2^{16}=65\,536$ 个盘块，分区最大 $=65\,536\times 4\text{ KiB}=256\text{ MiB}$。顺带可得 FAT 表本身占 $65\,536\times 2\text{ B}=128\text{ KiB}$。

\boxed{\text{答案：（1）4 MiB；（2）256 MiB}}\else \answerarea[5] \fi

\vspace{8pt}

\item \textbf{（单项选择题）} 某文件 \texttt{/home/u/a.txt} 被用户创建了两个硬链接和一个软链接，下列说法正确的是
\begin{center}
\begin{tabular}{p{0.47\linewidth}p{0.47\linewidth}}
(A) 删除其中一个硬链接后，文件数据立即被释放 & (B) 软链接可以跨文件系统，硬链接不可以 \\[4pt]
(C) 硬链接可以指向目录，软链接不可以指向目录 & (D) 每建立一个硬链接，系统就为文件新建一个 i-node
\end{tabular}
\end{center}

\ifshowsolutions\par\noindent\textbf{解析：} 硬链接只是让一个新目录项指向已有的 i-node，i-node 中的链接计数加 1，并不复制数据、也不新建 i-node，故 (D) 错。删除一个硬链接只是计数减 1，计数为 0 时（本例中两个硬链接都删掉）才释放 i-node 和数据块，故 (A) 错。

硬链接的目录项里存的是 i-node 号，而 i-node 号只在同一个文件系统内有意义，所以硬链接不能跨文件系统；软链接是一个独立文件，内容为目标路径字符串，由路径解析去定位目标，因此可以跨文件系统，(B) 正确。

(C) 说反了：为防止目录树出现环，通常禁止对目录建硬链接，而软链接指向目录是允许的。\boxed{\textbf{答案：}(B)}\else \answerarea \fi

\end{enumerate}

\subsection{真题风格与综合训练}

\begin{enumerate}[start=7]

\item \textbf{（单项选择题）} 某文件系统的 i-node 含 10 个直接块指针、1 个一级间接块指针、1 个二级间接块指针和 1 个三级间接块指针。盘块大小 4 KiB，每个盘块号占 4 字节。某文件已被打开（其 i-node 已在内存），现要读取该文件内偏移量 $5\text{ MiB}$ 处的一个字节，且相关的间接块与数据块均不在内存。该字节所在区域及为此需要的磁盘访问次数是
\begin{center}
\begin{tabular}{ll}
(A) 一级间接区，2 次 & (B) 二级间接区，2 次 \\[4pt]
(C) 二级间接区，3 次 & (D) 三级间接区，4 次
\end{tabular}
\end{center}

\ifshowsolutions\par\noindent\textbf{解析：} 思路：先把字节偏移换算成逻辑块号，再按“直接块 $\to$ 一级 $\to$ 二级 $\to$ 三级”的顺序逐段扣除，看落在哪一段；访问次数则数“从 i-node 出发到数据块之间要读几个块”。

每个索引块可存 $4\text{ KiB}/4\text{ B}=1024$ 个块号。偏移 $5\text{ MiB}=5\times 1024\text{ KiB}$ 对应的逻辑块号为 $5\times 1024\text{ KiB}/4\text{ KiB}=1280$。

容量分段：直接块覆盖逻辑块 $0\sim 9$；一级间接块覆盖 $1024$ 块，即逻辑块 $10\sim 1033$；二级间接块从逻辑块 $1034$ 起。$1280\ge 1034$，所以该字节在二级间接区。

访问路径：i-node 的二级间接指针 $\to$ 二级间接块 $\to$ 该项指向的一级间接块 $\to$ 数据块。i-node 已在内存，不计读盘；三个块都不在内存，故依次读二级间接块、一级间接块、数据块，共 3 次磁盘访问。\boxed{\textbf{答案：}(C)}\else \answerarea \fi

\vspace{8pt}

\item \textbf{（单项选择题）} 某系统用位示图管理磁盘空闲空间，字长 32 位，“1 表示已分配、0 表示空闲”，第 $i$ 个字的第 $j$ 位（$j$ 自低位向高位计）对应盘块号 $i\times 32+j$。分配时按字号递增、同一字内位号递增的顺序找第一个 0 位。已知所有字号小于 100 的字均已没有 0 位，而字号 100 的那个字的值为 \texttt{0xFFFF0000}，则本次分配的盘块号是
\begin{center}
\begin{tabular}{ll}
(A) 3200 & (B) 3215 \\[4pt]
(C) 3216 & (D) 3231
\end{tabular}
\end{center}

\ifshowsolutions\par\noindent\textbf{解析：} 思路：\texttt{0xFFFF0000} 是一个 32 位十六进制数，高 16 位（第 16～31 位）全为 0，低 16 位（第 0～15 位）全为 1。位号按“自低位向高位”递增，说明低号位先被检查，而低号位恰好全部已分配，所以第一个空闲位是 $j=16$。

字号 100 管理的盘块号范围是 $100\times 32=3200$ 到 $3231$，对应的位号 $j=0\sim 31$。代入 $b=i\times 32+j=100\times 32+16=3216$。

常见错误是把最高位当作位号 0，那样会算出 $j=0$ 从而选 3200；也有人把“第一个 0 位”误当作最后一个 1 位而选 3215，两者都不符合题设的编号约定。\boxed{\textbf{答案：}(C)}\else \answerarea \fi

\vspace{8pt}

\item \textbf{（真题风格·综合应用题）} 某 UNIX 类文件系统的 i-node 含 12 个直接块指针、1 个一级间接块指针、1 个二级间接块指针和 1 个三级间接块指针。盘块大小为 1 KiB，每个盘块号占 4 字节，一个指针恰好指向一个盘块。
（1）该文件系统单个文件的最大长度是多少字节？用 GiB/MiB/KiB 的形式和字节数各表示一次。
（2）文件已打开、i-node 在内存，要读取文件内偏移量 $20\text{ MiB}$ 处的一个字节，该字节落在哪一级间接区？若相关的间接块与数据块都不在内存，需要几次磁盘访问？该字节对应的块号在二级间接块中处于什么位置？
（3）若同一磁盘改用显式链接（FAT）管理，盘块仍为 1 KiB，FAT 表项 16 位，则最大分区容量和 FAT 表本身的大小各是多少？

\ifshowsolutions\par\noindent\textbf{解析：} 思路：先算“一个索引块能装多少个块号”，这是所有小问的公共常数；再按容量分层累加；访问次数则沿“i-node $\to$ 各级索引块 $\to$ 数据块”数经过的块。

块号 4 字节、块 1 KiB，故每个索引块可存 $1024/4=256$ 个块号。

\textbf{（1）} 分层计算：
\begin{align*}
\text{直接块：}\quad & 12\times 1\text{ KiB}=12\text{ KiB}\\
\text{一级间接：}\quad & 256\times 1\text{ KiB}=256\text{ KiB}\\
\text{二级间接：}\quad & 256^{2}\times 1\text{ KiB}=65\,536\text{ KiB}=64\text{ MiB}\\
\text{三级间接：}\quad & 256^{3}\times 1\text{ KiB}=16\,777\,216\text{ KiB}=16\text{ GiB}
\end{align*}
最大长度 $=16\text{ GiB}+64\text{ MiB}+268\text{ KiB}$。换成字节：$16\times 2^{30}=17\,179\,869\,184$，$64\times 2^{20}=67\,108\,864$，$268\times 2^{10}=274\,432$，三项相加得 $17\,247\,252\,480$ 字节。

\textbf{（2）} 逻辑块号 $=20\text{ MiB}/1\text{ KiB}=20\,480$。直接块覆盖 $0\sim 11$；一级间接覆盖 $12\sim 12+255=267$；二级间接覆盖 $268\sim 268+65\,535=66\,803$。$20\,480$ 落在二级间接区。

在二级间接区内的序号为 $20\,480-268=20\,212$。二级间接块有 256 项，$20\,212=78\times 256+244$，所以它是二级间接块第 78 项所指向的那个一级间接块中的第 244 项（均从 0 编号，第 78 项指向的一级块负责二级区的第 $78\times 256\sim 78\times 256+255$ 块）。

读盘次数：i-node 已在内存，先读二级间接块（1 次），再读它第 78 项指向的一级间接块（1 次），最后读一级块第 244 项指向的数据块（1 次），共 3 次。

\textbf{（3）} 16 位表项可表示 $2^{16}=65\,536$ 个盘块，分区最大 $65\,536\times 1\text{ KiB}=64\text{ MiB}$；FAT 每块一项、每项 2 字节，占用 $65\,536\times 2\text{ B}=128\text{ KiB}$。

\boxed{\text{答案：（1）$17\,247\,252\,480$ 字节；（2）3 次；（3）64 MiB 与 128 KiB}}\else \answerarea[12] \fi

\vspace{8pt}

\item \textbf{（真题风格·综合应用题）} 某磁盘有 200 个柱面，编号 $0\sim 199$，磁头当前位于 53 号柱面，并正朝柱面号增大的方向移动。柱面请求队列（按到达顺序）为
\[
98,\ 183,\ 37,\ 122,\ 14,\ 124,\ 65,\ 67 .
\]
忽略旋转延迟与传输时间，只统计寻道距离。约定：SCAN 沿当前方向一直服务到物理端点 199 后才折返；C-SCAN 只在一个方向服务，到达 199 后立即移回 0 号柱面（这段返回移动计入寻道距离），再从 0 开始只向增大方向服务。
（1）按 FCFS 服务，写出服务顺序并求总寻道距离与平均寻道距离。
（2）按 SSTF 服务，写出服务顺序并求总寻道距离与平均寻道距离。
（3）按 SCAN 和 C-SCAN 服务，分别写出服务顺序并求总寻道距离与平均寻道距离。
（4）就本题数据说明四种算法在总寻道距离和公平性上的差别。

\ifshowsolutions\par\noindent\textbf{解析：} 思路：总寻道距离就是“相邻两次服务位置之差的绝对值”之和，起点的第一次移动从 53 算起；平均寻道距离 $=$ 总距离 $/$ 请求数 $=8$。SCAN 与 C-SCAN 的关键区别在于是否把返回段计入，以及返回段是否服务请求。

\textbf{（1）FCFS} 按到达顺序：$53\to98\to183\to37\to122\to14\to124\to65\to67$。
\[
45+85+146+85+108+110+59+2=640
\]
平均 $640/8=80$ 个柱面。

\textbf{（2）SSTF} 每次选离当前位置最近的请求：
\[
53\to65\to67\to37\to14\to98\to122\to124\to183
\]
各段距离 $12,2,30,23,84,24,2,59$，总和 $236$，平均 $236/8=29.5$ 个柱面。注意 67 之后最近的未服务请求是 37（距离 30），而不是 98（距离 31）。

\textbf{（3）SCAN}（方向增大，到 199 折返）：$53\to65\to67\to98\to122\to124\to183\to199\to37\to14$。
\[
(12+2+31+24+2+59+16)+(199-37)+(37-14)=146+162+23=331
\]
其中 $53\to199$ 段共 146，$199\to37$ 段 162，$37\to14$ 段 23。平均 $331/8=41.375$ 个柱面。

\textbf{C-SCAN}：$53\to199\to0\to14\to37$。
\[
(199-53)+199+(14-0)+(37-14)=146+199+14+23=382
\]
平均 $382/8=47.75$ 个柱面。若约定返回段不计入寻道，则为 $146+14+23=183$；本题按题设口径计入，故取 382。

\textbf{（4）} 总寻道距离：SSTF（236）$<$ SCAN（331）$<$ C-SCAN（382）$<$ FCFS（640），本题中 SSTF 最优。但 SSTF 只顾“最近”，若新请求不断到达磁头附近，远端请求可能长期得不到服务（饥饿）；SCAN、C-SCAN 有确定的扫描方向，不会饥饿，其中 C-SCAN 把等待时间摊得更均匀，代价是返回段的空走。

\boxed{\text{答案：（1）640；（2）236；（3）SCAN 331、C-SCAN 382；（4）SSTF 最短但可能饥饿}}\else \answerarea[12] \fi

\end{enumerate}
